% ============================================================================
% CHAPTER 4 - PROPOSED METHODOLOGY   (template file chapter_5.tex)
% Rubric mapping: CO5 (15 marks, 5 supervisor + 10 defense panel) is assessed
% from 4.1 Design Process and 4.2 Preliminary Design.
% ============================================================================

\section{Design Process or Methodology Overview}

\subsection{What the design has to achieve}

The design follows directly from the objectives in Section 1.4 and the constraints
in Section 3.1.5. A lecture video goes in. A lecture note comes out that contains
the lecturer's words and the lecturer's board. The whole thing has to run on one
desktop graphics card, it has to be measurable stage by stage, and it must not
invent content.

\subsection{Why a pipeline of stages rather than one model}

One design option was a single multimodal model that reads the video and writes
the note. It was rejected for three reasons.

The first is measurement. A single model gives one number at the end, and when
that number is poor there is nothing to look at. A pipeline gives a separate
measurement for the speech, for the board reconstruction, for the board reading
and for the note, so a weakness can be located. Chapter 5 shows this paying off,
because the four stages turned out to contribute very differently and one of them
contributes almost nothing to the metric used.

The second is hardware. Keeping the stages apart means only one model is in memory
at a time. The speech model, the vision-language model and the note-writing model
never have to fit together.

The third is honesty about content. When the board is read into text as a separate
step, the note-writing model receives the board as text rather than having to
remember it, so it has less reason to invent. When the note is written by a model
that also had to look at the image, there is no way to tell whether a fact came
from the board or from the model's training data.

\subsection{The four stages}

\figplaceholder{fig-4-1-pipeline}{The InsightLens pipeline. Each stage writes a
file that the next stage reads.}{The main system diagram, drawn left to right or
top to bottom. Boxes in order: (1) Video, (2) Audio extraction at 16 kHz and
frame extraction every 2 s, (3) Stage 1 Speech: Whisper large-v3-turbo plus LoRA
adapter, output "transcript.txt (timestamped Banglish)", (4) Stage 2 Board:
erase detection, era split, tiled mosaic, person mask, clean-up, output "one
clean board image per era", (5) Stage 3 Board reading: box finding, numbered
coloured boxes drawn, Qwen2.5-VL names and transcribes each box, output
"board\_boxes.json", (6) Stage 4 Notes: Qwen2.5-7B writes one section per board,
quote checker, box-reference check, output "notes in Banglish and English,
Markdown and HTML". Show the file passed between each pair of stages on the
arrow. Mark which stages run on which machine if there is room.}{9cm}

\textbf{Stage 1, speech.} The audio is extracted and cut into clips. A LoRA
adapter trained on the corpus is attached to Whisper large-v3-turbo, and the
adapted model writes a timestamped transcript in romanized Banglish.

\textbf{Stage 2, the board.} Frames are extracted at a fixed interval. Erase
events are detected and split the lecture into eras, one board per era. Within an
era the frame is tiled, and each tile is taken from a moment when the lecturer was
not standing in front of that tile. The result is one image per board in which
every pixel came from some real frame.

\textbf{Stage 3, reading the board.} Blocks of writing are found on the clean
board, a numbered coloured box is drawn around each, and Qwen2.5-VL is asked to
name and transcribe each numbered box.

\textbf{Stage 4, the note.} For each board, Qwen2.5-7B-Instruct is given the boxes,
the board text and the part of the transcript that belongs to that board, and
writes one section of notes that points at the boxes. Checkers then remove
unverifiable quotations and bad box references, and the sections are assembled
into a page.

\subsection{Principles that were fixed before building}

Four rules were adopted at the start and were not relaxed.

\textbf{Every stage is a separate program with a file interface.} A stage is
skipped when its output already exists, so a run continues after an interruption
and a single stage can be replaced without touching the others.

\textbf{Nothing is invented in the image.} The board reconstruction may leave a
visible gap, but it may not fill one with plausible pixels.

\textbf{The answer keys are never given to a generator.} The board answer keys
used for scoring sit in a folder that the note generator refuses to read, enforced
in code rather than by discipline.

\textbf{A number needs a command.} Every result in Chapter 5 is produced by a
command recorded beside it.

\section{Preliminary Design and Alternatives Considered}

Several designs were built and measured before the final one was chosen. This
section presents them with the verification each received, because three of them
failed and the failures are part of the argument for the design that was kept.

\subsection{Alternatives for the speech stage}

\textbf{Option A, use Whisper as it is.} This is the obvious baseline and it was
measured first. On the held-out lectures it gives a median character error rate of
67.7 per cent and a median word error rate of 93.9 per cent, and it translates the
speech into English rather than transcribing it. Rejected on the measurement.

\textbf{Option B, use a Bengali speech model and transliterate.} A Bengali model
writes the Bengali script, so the output would have to be romanized afterwards.
This loses the English technical terms twice, first when the Bengali model spells
them phonetically in Bengali and again when they are romanized back. Rejected on
the design, without a full experiment, because the failure is structural.

\textbf{Option C, run two recognisers and merge them.} An English recogniser and a
Bengali recogniser could be run on the same audio and their outputs combined.
This was built and measured on nine lectures. The merged output scored 0.7
percentage points above the baseline on the term measure, which a paired
permutation test puts at p = 0.32, and its word error rate was far worse than the
baseline at 146.8 per cent against 82.6 per cent. Rejected on the measurement.

\textbf{Option D, bias the decoder towards words seen on the board.} The idea was
to push the speech model towards terms the board shows, implemented as a logits
processor with a strength parameter. A sweep over the strength from 0 to 2.0 found
that any non-zero value halves term recall and truncates the transcript, and that
every non-zero value produces byte-identical output. Rejected on the measurement,
and the mechanism was disabled.

\textbf{Option E, adapt Whisper to this speech with LoRA.} Selected. It keeps the
multilingual model's acoustic knowledge, trains a small number of parameters so
that a few hours of speech is enough, and produces the romanized output directly
because that is what the training references contain.

\subsection{Alternatives for reconstructing the board}

\textbf{Option A, take the clearest single frame of each era.} Simple, and it is
the fallback the pipeline still supports. It fails whenever the lecturer stands in
front of part of the board in every frame of the era, which measurement showed to
be the normal case rather than the exception. Over 79 frames of one lecture the
lecturer covers 23.5 per cent of the writing surface in the median frame and never
less than 5 per cent.

\textbf{Option B, remove the lecturer and fill the hole by inpainting.} This was
the method in the first version of the project. Inspection of its output against
the source frame found two failures. The lecturer remained visible as a grey
silhouette because the brightness threshold missed the light squares of his
shirt, and text he was writing at that moment was lost because the filled region
was rebuilt from a single donor frame. Worse, the filled figure showed content
that had not yet been written at that time, so the image was a plausible picture
of a board that never existed. Rejected.

\textbf{Option C, generative inpainting with a diffusion model.} Not attempted, on
purpose. On a whiteboard, a plausible fill is invented writing, and a reader
cannot tell it from the real thing. A visible gap is honest and an invented
character is not.

\textbf{Option D, a tiled mosaic over erase-separated eras.} Selected, and
described in Section 4.4.3. Its property is that every output pixel is unmodified
camera output, so there is no fidelity question to argue about. The only question
is whether each tile found a clear moment, and that is counted exactly and
reported.

\subsection{Alternatives for reading the board}

\textbf{Option A, classical optical character recognition.} Standard OCR engines
are built for printed text. The board carries handwriting, diagrams with arrows,
truth tables and code, and the earlier phase of the project already showed this
route giving very little.

\textbf{Option B, ask the vision-language model for keywords.} This is what the
first version of the pipeline did. Measured against the hand-verified answer keys
over 35 boards it recovers 31.2 per cent of the items and, tellingly, none of the
67 numbers. Rejected on the measurement.

\textbf{Option C, ask the vision-language model to transcribe the board in full.}
Selected. The same model with the same images recovers 88.8 per cent of the items
and 65 of the 67 numbers. The prompt asks for every number exactly as written and
for the marker \texttt{[illegible]} instead of a guess.

\textbf{Option D, let the model find the regions itself.} The code supports asking
the model to draw the boxes, but it was not adopted, because our own box finder is
deterministic and its output can be drawn onto the real image and checked by eye.
This option remains untested and is listed as such.

\subsection{Alternatives for writing the note}

\textbf{Option A, ask for good notes on the lecture.} The original prompt. The
model recited textbook material, including a definition of a NAND gate that the
lecturer never gave and that was wrong. Measured board recall 37.2 per cent over
35 boards.

\textbf{Option B, a grounded prompt with board keywords as the visual input.}
Recall fell to 25.8 per cent. The grounded prompt stops the model reciting, and
with only keywords there is nothing better to put in place of the recitation.

\textbf{Option C, a grounded prompt with the full board transcription.} Recall
88.0 per cent. This works, but the note reads like a paste of the board text.

\textbf{Option D, the annotated note.} Selected. One section per board, pointing
at the numbered boxes, with the lecturer quoted, a step by step walk through the
board, and non-lecture material confined to a labelled box. It reaches 89.1 per
cent, which is statistically indistinguishable from option C, while being a
readable document rather than a paste.

For the English version two routes were built and compared. Writing English
directly from the board and the transcript reaches 55.9 per cent on this measure.
Writing the section in Banglish first and then translating it reaches 89.4 per
cent. The translated route was selected, and Chapter 5 explains why the gap is
partly an artefact of the measure and reports the caveat that four dense boards
carry most of it.

\subsection{Summary of the design decisions}

\begin{table}[H]
\centering
\caption{The alternatives considered at each stage and the reason for the choice.
Options marked as measured were built and evaluated before being rejected.}
\label{tab:alternatives}
\small
\begin{tabular}{p{2.1cm}p{4.6cm}p{2.0cm}p{5.0cm}}
\toprule
\textbf{Stage} & \textbf{Option} & \textbf{Status} & \textbf{Reason} \\
\midrule
Speech & Off-the-shelf Whisper & measured &
CER 67.7 per cent, and it translates \\
Speech & Bengali model plus transliteration & rejected by design &
loses English terms twice \\
Speech & Dual recogniser fusion & measured &
+0.7 pp, p = 0.32, and WER far worse \\
Speech & Visual bias in the decoder & measured &
halves term recall at any strength \\
Speech & \textbf{Whisper plus LoRA} & \textbf{selected} &
CER 15.8 per cent on unseen lectures \\
\addlinespace
Board & Best single frame & fallback &
the lecturer is never absent \\
Board & Mask and inpaint & measured &
left a silhouette and showed writing from the wrong time \\
Board & Generative inpainting & not attempted &
would invent writing \\
Board & \textbf{Tiled mosaic over eras} & \textbf{selected} &
every pixel is real, coverage is counted \\
\addlinespace
Reading & Classical OCR & rejected &
built for print, not handwriting and diagrams \\
Reading & Keyword prompt & measured &
31.2 per cent recall, no numbers at all \\
Reading & \textbf{Full transcription prompt} & \textbf{selected} &
88.8 per cent recall \\
Reading & Model draws its own boxes & untested &
kept in the code, not evaluated \\
\addlinespace
Notes & Free summarisation & measured & 37.2 per cent, recites a textbook \\
Notes & Grounded prompt, keywords & measured & 25.8 per cent \\
Notes & Grounded prompt, board text & measured & 88.0 per cent, reads like a paste \\
Notes & \textbf{Annotated note with boxes} & \textbf{selected} &
89.1 per cent and readable \\
\bottomrule
\end{tabular}
\end{table}

\section{Data Collection}

\subsection{The recordings}

Forty-four lectures were recorded in Brac University classrooms with a fixed
camera pointed at the lecturer and the whiteboard. Three lecturers took part, and
no lecture uses slides. The subjects are introductory programming, digital logic,
databases and computer networking. The recordings total 8.33 hours, with the
shortest at 3.4 minutes and the longest at 36.4 minutes.

The lecturer of each recording was confirmed by voice rather than by memory. A
WavLM speaker embedding was computed for every lecture and compared with reference
embeddings, giving cosine similarity between 0.98 and 0.998 within a lecturer and
between 0.59 and 0.90 across lecturers. This check exists because an early version
of the split had one lecture attributed to the wrong lecturer, which made a result
look better than it was, and the error is described in Section 
ef{sec:asr-loso}.

\subsection{The transcription standard}

Before any transcription started, a written standard was produced, and every
transcript is checked against it by a validation script. Its main rules are as
follows.

\begin{itemize}
  \item Each line is a timed segment written as \texttt{[M:SS-M:SS]} followed by
        the text, with no spaces inside the brackets.
  \item Segments are 10 to 25 seconds long and never longer than 30, because
        Whisper processes a 30 second window.
  \item There are no gaps between segments and no overlaps.
  \item Fillers are kept, including \textit{aaa}, \textit{um}, \textit{mane},
        \textit{je} and \textit{tahole}, because the goal is a transcript of the
        speech and not a tidied version of it.
  \item The text is ASCII only, with no smart quotes.
  \item A table of canonical spellings fixes the most common words, so that the
        same word is not written three ways by three transcribers.
  \item A header records the lecture, the video file, its duration and the
        lecturer.
\end{itemize}

Transcription was done by hand. Every word of every transcript was checked against
the audio by a member of the team. This costs about one hour of work for ten
minutes of video.

\subsection{Data Cleaning}

Four cleaning problems occurred and each is handled by a script rather than by
hand editing.

\textbf{Malformed timestamps.} The segment parser accepts a timestamp only in the
exact form \texttt{[1:34-1:53]}. In an earlier version of the corpus, 9 of 156
timestamps carried stray spaces, written for example as \texttt{[4:22- 4:45 ]}.
Those lines did not register as segment boundaries, so their text was joined to the
previous segment and ten evaluation clips carried more words than their audio. The
validator now reports every unreadable timestamp and can repair them with a
\texttt{--fix} flag, keeping a backup copy of the original file. The effect of the
defect on the affected result was measured rather than assumed, and is reported in
Section 
ef{sec:asr-loso}.

\textbf{Files not yet ready.} A transcript that is still being written is marked
by its file name, and the pipeline skips it until the marker is removed. Files
were also delivered at one point with a placeholder extension that made them
invisible to the pipeline, which was found by counting how many transcripts the
loader actually saw rather than how many files existed.

\textbf{Missing lecturer labels.} Where the lecturer header was missing, it was
filled from the voice check rather than guessed, and any recording where the voice
check was not confident was left for a person to decide.

\textbf{Duplicate or mismatched names.} A preflight script refuses to start
training if a transcript has no matching video, if two different videos share a
name, if a transcript is longer than its video, or if a lecturer label disagrees
with the voice check. This runs as a hard gate before every training run, because
all of these faults corrupt training silently rather than crashing.

\subsection{Data Transformation}

\textbf{Audio.} Each video is converted to 16 kHz mono PCM WAV, which is the input
format Whisper expects.

\textbf{Clips.} The human transcript gives the segment boundaries, so a clip is cut
for each segment. Segments longer than the model's 30 second window are subdivided
at sentence ends, with a working limit of 28 seconds to stay comfortably inside the
window. Each clip becomes one training example of an audio array and its reference
text.

\textbf{Features.} The audio is turned into a log-Mel spectrogram by the model's own
feature extractor, and the reference text is tokenised by the model's tokeniser, so
that no custom preprocessing can differ between training and evaluation.

\textbf{Frames.} For the vision side, frames are extracted at a fixed interval and
scaled so that the longer side is at most 1920 pixels. A 10 second interval is
enough for lecturers who move about. For a lecturer who stands still, a 2 second
interval gives each tile five times as many chances to see the board clear, and
Section 
ef{sec:board-coverage} shows what that is worth.

\subsection{Data Integration}

The corpus was renumbered part way through the project so that lectures are grouped
by lecturer. Every result produced before the renumbering used the old names, so
two things were done rather than renaming files and hoping. A lookup table records
the correspondence between old and new numbers, and every script that has to find a
lecture's audio or board uses the table instead of assuming that the number in a
file name means the same thing everywhere. The exact ground truth behind the earlier
results is frozen in a separate folder, so any of those results can still be
reproduced by pointing the data preparation script at the frozen copy.

Three sources have to line up for one lecture, which are the video, the transcript
and the reconstructed boards. They are joined by lecture number through that table,
and the preflight script checks the join before any training.

\subsection{Data Reduction}

No dimensionality reduction is applied to the audio or the images, since the models
consume raw signals. Reduction happens in three other places.

Lectures without a transcript take no part in training or in the speech evaluation,
which leaves 28 of the 44 for the speech experiments. All 44 are available to the
vision side, because the vision stages never touch the audio.

For scoring the notes, only the 13 lectures with hand-verified board answer keys can
be used. The other 30 lectures receive the same notes as demonstrations and are not
scored, and the report says which is which every time.

Within the board stage, an era shorter than four frames is ignored, because too few
frames means too few chances for a tile to see the board clear and the resulting
board would be poor evidence of anything.

\subsection{The splits}

Three separate splits are used and they must not be confused with one another.

\textbf{The validation split, for choosing settings.} Three lectures, one per
lecturer, were drawn with a fixed seed before any tuning started, and were recorded
in a file that the data preparation script reads. They are BanglaASR2 from lecturer
A, BanglaASR12 from B and BanglaASR14 from C, 34.2 minutes in total. They are used
only to score tuning runs and are locked out of every test set afterwards by a
\texttt{--never-test} flag.

\textbf{The test split, for the final result.} Whole lectures are held out at
random, about 20 per cent of each lecturer's minutes, with the seed fixed at zero
and the rule written down before the draw. A lecture is never cut between training
and test, and every lecturer appears on both sides. The draw selected BanglaASR8, 9,
11, 15, 19 and 27, which is 1.05 hours and 177 clips, leaving 4.10 hours and 22
lectures for training. Because the test lecturers are also heard in training, this
design answers the question ``new lectures from known lecturers'', and the report
labels it that way every time it is quoted.

\textbf{The leave-one-speaker-out splits, for the harder question.} Each lecturer in
turn is held out entirely and the model is trained on the other two. This answers
``an unseen lecturer'', which is harder and gives worse numbers, and both are
reported.

There is one more rule, which exists to protect the note evaluation. A note that is
scored on a lecture's board answer key must be written from a transcript produced by
an adapter that never heard that lecture's lecturer. The transcripts used for the
scored lectures therefore come from the leave-one-speaker-out adapters, and the
builder records in each output file which adapter produced its input.

\subsection{Summary of Preprocessed Data}

\begin{table}[H]
\centering
\caption{The corpus after preprocessing.}
\label{tab:corpus}
\small
\begin{tabular}{p{7.0cm}p{6.4cm}}
\toprule
\textbf{Quantity} & \textbf{Value} \\
\midrule
Lecture recordings & 44, 8.33 hours \\
Lectures with a hand-checked transcript & 28, 5.15 hours of timed speech \\
Lectures used for vision only & 16, 3.10 hours \\
Transcribed segments & 609 \\
Median segment length & 24 seconds \\
Words transcribed & about 39,100 \\
Lecturer A & 9 lectures, 1.51 h, 141 segments \\
Lecturer B & 4 lectures, 0.97 h, 39 segments \\
Lecturer C & 15 lectures, 2.66 h, 429 segments \\
Training set for the final model & 22 lectures, 4.10 hours \\
Test set for the final model & 6 lectures, 1.05 hours, 177 clips \\
Validation lectures, never tested on & 3 lectures, 34.2 minutes \\
Reconstructed boards & 145 across 40 lectures \\
Boards with a hand-verified answer key & 45, holding 436 items \\
Boards checked by hand for missing writing & 98 \\
\bottomrule
\end{tabular}
\end{table}

Figure \ref{fig:fig-data-composition} shows how the recordings divide between
lecturers and between training, testing and vision-only use. The imbalance is real
and matters for the results, since lecturer B contributed the least speech and
Chapter 5 shows that lecturer B's test lecture improves the least.

\begin{figure}[H]
\centering
\includegraphics[width=0.85\textwidth]{fig-data-composition}
\caption{The corpus by lecturer. The darker part of each bar is training speech,
the middle part is held out for testing, and the lightest part is video with no
transcript, used for the vision side only.}
\label{fig:fig-data-composition}
\end{figure}

\begin{figure}[H]
\centering
\includegraphics[width=0.85\textwidth]{fig-data-segments}
\caption{Lengths of the 609 transcribed segments against Whisper's 30 second
window. Segments longer than the window are subdivided at sentence ends before
training.}
\label{fig:fig-data-segments}
\end{figure}

\begin{figure}[H]
\centering
\includegraphics[width=\textwidth]{fig-data-per-lecture}
\caption{Transcribed minutes and speaking rate for each lecture with a transcript.
Bar colour is the lecturer. Lectures marked \textit{test} are the six held out for
the final result.}
\label{fig:fig-data-per-lecture}
\end{figure}

\begin{figure}[H]
\centering
\includegraphics[width=\textwidth]{fig-data-boards}
\caption{Left, the number of erase-separated boards found in each lecture. Right,
the distribution of the fraction of tiles that found a completely clear view of
the board, with the median marked.}
\label{fig:fig-data-boards}
\end{figure}

\section{Implementation of Selected Design}

\subsection{Stage 1: adapting the speech model}

The base model is whisper-large-v3-turbo, an encoder-decoder transformer with 809
million parameters \cite{radford2023whisper}. The base weights are frozen. LoRA
adapters are attached to the query and value projection matrices of the attention
layers, with rank 16 and scaling factor alpha 32 \cite{hu2022lora}. Training uses a
batch size of 8, 8 epochs, bfloat16 precision, and gradient checkpointing on the
12 GB card, which cuts the memory needed from about 18 GB to 6.1 GB at the cost of
recomputing activations. Every run is done with two random seeds, 42 and 1, and both
are reported.

\textbf{Decoding.} The reported numbers use plain greedy decoding. The loop
safeguard from the Whisper paper is also available and is reported beside the main
result for comparison. It works by checking the gzip compression ratio of a
hypothesis, and if the ratio is above 2.4 the clip is decoded again with sampling at
temperatures 0.2, 0.4 and upwards until the output stops looping. It reads only the
hypothesis and never the reference, and it is applied identically to the base and
the adapted model. Only the compression ratio half of the published rule is used,
because the log probability half would resample most Banglish clips rather than only
the looping ones.

\subsection{Stage 1b: choosing the settings, and the pre-registered plan}

The panel question that any tuned result invites is whether the settings were
chosen by trying things until the test number looked good. The answer here is a
written plan, committed to version control before any tuning run, that fixes the
validation lectures, the order of the search, and the rule for accepting a winner.

The plan has five stages, run in order, each one starting from the winner of the
previous stage.

\begin{enumerate}
  \item \textbf{Learning rate}, over 5e-4, 1e-3 and 2e-3.
  \item \textbf{LoRA rank}, over 8, 16 and 32, with alpha fixed at twice the rank.
  \item \textbf{Which layers are adapted}, comparing the query and value
        projections against all four attention projections.
  \item \textbf{Number of epochs}, taken from the validation loss curve.
  \item \textbf{Seed stability}, by repeating the winner with a second seed.
\end{enumerate}

Everything else is held fixed, which means the base model, batch size 8, bfloat16,
gradient checkpointing and greedy decoding. The score is the median per-clip
character error rate on the validation clips. The acceptance rule is that a
candidate replaces the default only if it beats the default by more than the spread
between two seeds of the same configuration, which is the plan's way of refusing to
chase noise.

The procedure was run twice. A rehearsal on 2.69 hours of transcript was used to
check that the machinery worked, and the real run was repeated on the full corpus
once transcription closed. The second run supersedes the first, and Section 
ef{sec:asr-tuning}
reports both.

\subsection{Stage 2: reconstructing the board}

\textbf{Finding the eras.} The amount of ink on the board is tracked across frames.
When 30 per cent of it disappears between consecutive frames, the board has been
wiped, and the frames between two wipes form one era. An era shorter than four
frames is discarded.

\textbf{Tiling.} Work is done at a scan width of 640 pixels. The frame is cut into
tiles of 32 scan pixels with 50 per cent overlap between neighbours. A tile counts
as clear in a given frame when less than 2 per cent of it is covered by the
lecturer. For each tile independently, the pipeline takes the latest frame within
the era in which that tile was clear, and the chosen tiles are blended back together
with a raised-cosine weight so that the seams do not show.

This works because the lecturer occupies one place at a time. It is also the reason
the method has a stated limit, which is that it needs the lecturer to move. Where a
tile never finds a clear moment, the residue is left visible and reported rather
than filled in.

\textbf{Locating the lecturer.} Three options exist and the third is used. The
simplest takes whatever differs from the usual appearance of the board, which fails
when the lecturer stands still, because then he is part of the usual appearance, and
fails again when the camera's automatic exposure darkens a frame, because then every
pixel differs. The second uses a pretrained DeepLabV3 network with a ResNet-101
backbone and takes its person class frame by frame \cite{chen2017deeplabv3}. The
third adds a shadow mask to the network, computed as board that is darker than usual
after dividing out each frame's exposure gain, with pen-stroke-thin marks eroded away
so that writing is not mistaken for shadow. The third option recovers writing the
first one lost, including a diagram drawn at the end of an era and seven of the eight
digits of a binary number that the old mask reduced to four.

\textbf{Display clean-up.} A separate step flattens the background and fades
everything that is not clearly darker than its surroundings to white, in the manner
of the whiteboard mode of a document scanner application, and blanks the pixels the
mosaic itself flagged as the lecturer. This makes every board a clean white image.
It is important to be exact about what this step is, because it replaces background
pixels. A cleaned board is therefore not unmodified camera output. The strokes are
camera pixels and the background is not, and the report says so wherever a cleaned
board is shown. Section 
ef{sec:board-cleanup} reports the measured effect of this step on how well
the vision-language model reads the board, which was not the expected one.

\figplaceholder{fig-4-2-before-frame}{A single frame of the same lecture, showing
why one frame is not enough. The lecturer stands in front of the part of the board
he has just written on.}{Take one frame from
\texttt{output/lectures/BanglaASR29/frames/} where the lecturer is standing in
front of the writing, for example \texttt{frame\_000619\_001238s.jpg}, and blur or
pixelate his face before using it. The ethics statement says no person is
identifiable in the figures, so the face must be covered. Put it beside or above
the reconstructed board in Figure \ref{fig:board-mosaic-raw} if you prefer a
two-panel figure.}{6.5cm}

\begin{figure}[H]
\centering
\includegraphics[width=0.92\textwidth]{board-mosaic-raw}
\caption{A reconstructed board from lecture BanglaASR29, era 5. Every pixel comes
from a real frame of the video. The lecturer stood in front of the left half of
this board for most of the era, and what remains of him is the soft grey shadow on
the left, not a filled-in guess.}
\label{fig:board-mosaic-raw}
\end{figure}

\begin{figure}[H]
\centering
\includegraphics[width=0.92\textwidth]{board-mosaic-clean}
\caption{The same board after the display clean-up. The background has been
replaced, so this image is enhanced rather than unmodified camera output, while
the strokes themselves are camera pixels.}
\label{fig:board-mosaic-clean}
\end{figure}

\subsection{Stage 3: numbered boxes and the vision-language model}

\textbf{Finding the boxes.} Our own code, not the model, finds the blocks of
writing. It groups dark pixels, strips the board frame, drops specks by counting
dark pixels, and rejoins boxes that are parts of the same written line. At most ten
boxes are produced per board. Each box gets a number and a colour from a fixed list
of ten distinct colours, and the notes later refer to a box by both, for example
``the purple box 5''.

\textbf{Naming and transcribing.} The numbered boxes are drawn onto the clean board
and the image is given to Qwen2.5-VL-7B-Instruct with a prompt that asks, for each
numbered box, for a short name and a full transcription of what is inside it. This
is Set-of-Mark prompting \cite{yang2023som}, with the marks supplied by our code.
Boxes the model reports as empty are dropped.

Two defects in this stage were found during the real runs and both are handled in
code. On one board the model read the content correctly, then tried to draw a
diagonal line of the diagram and emitted the same backslash line 512 times, giving
42,037 characters. The reader now collapses a run of identical lines and records how
many it dropped, which matters because that text is pasted into the note prompt.
Separately, boards carrying code make the model emit unescaped quotation marks that
break the JSON it is asked to return, and without a fallback such a board loses all
of its boxes. A salvage parser now reads the identifier, name and text fields one at
a time when the JSON cannot be parsed.

\begin{figure}[H]
\centering
\includegraphics[width=0.92\textwidth]{board-demo-5}
\caption{The same board with the numbered coloured boxes drawn on it and named by
the vision-language model. The notes refer to these numbers and colours. Box 1 is a
manufacturer's sticker on the whiteboard, which the box finder treats as content
because it is dark and separate. This is a known cosmetic defect and is discussed in
Section 
ef{sec:known-defects}.}
\label{fig:board-demo-5}
\end{figure}

\subsection{Stage 4: writing the note}

The note is written by Qwen2.5-7B-Instruct, one call per board. Each call receives
the board image's boxes with their names and transcribed text, and the part of the
transcript that falls between that board's first and last frame, so that the model
sees the words that were spoken while that board was on the wall. A final call
produces the title, the takeaways and the self-check questions for the whole
lecture.

The prompt asks for a section that walks the reader through the board step by step,
refers to the boxes by number and colour, and quotes the lecturer between two and
five times. Material that is standard for the topic but that the lecturer did not say
is allowed, but only inside a section headed \texttt{Background}, which the page
builder renders as a dashed, labelled aside so that the reader can see it is not part
of the lecture. This is the controlled version of the failure that produced the wrong
textbook definition in the first version of the pipeline. The material is still
generated, but it is fenced and labelled instead of passing as something the lecturer
said.

Three checks then run over the generated text.

\textbf{The quote checker} takes every block quotation and searches for it in the
transcript. A quotation that is not there word for word is deleted, and the count of
kept and deleted quotations is written into the output file. A second counter flags
quoted spans of four or more words inside ordinary paragraphs that match neither the
transcript nor the board, but only flags them, because most of them turn out to be
code or translated lines rather than errors.

\textbf{The box reference check} removes any reference to a box number that the board
does not have, and counts them.

\textbf{The placeholder strip} removes the prompt's own example text when the model
copies it into the notes instead of leaving the quotation out. This was found by
reading the output, where the model had written the words \textit{exact words from
the transcript} as if they were a quotation.

\textbf{The two languages.} Each version is written wholly in its own language. The
Banglish note quotes the lecturer's real words, so the quote checker applies to it.
The English note quotes the lecturer in English translation, and a translation
cannot be matched against a Banglish transcript, so the checker does not run on the
English files. Every output file records which of the two it is, and Chapter 5 states
the resulting claim precisely, which is that quotations in the Banglish notes are
verified word for word and quotations in the English notes are unverified
translations.

The English note is produced by writing the section in Banglish first and then
translating it with the same model, keeping the Markdown structure, the tables, the
box numbers and the colours. A defect in this route was found and fixed. The
translation left 97 section labels in Banglish across the thirteen files, so a page
billed as English still showed Banglish headings, and three files returned the
summary line in capital letters. Both are repaired by a mapping script that does not
call the model, since the labels are a fixed set written by our own prompt. The
capitalisation repair takes each word's casing from how the same file writes that
word elsewhere, so that Python, TTL and DBMS survive.

\textbf{The page.} The sections are assembled into Markdown and into a single
self-contained HTML page with the board images embedded, coloured tags for the box
references and click-to-reveal answers for the self-check questions.

\subsection{How the system is run}

The whole pipeline is one command that takes a video and produces the notes. It runs
each stage as a separate process, skips any stage whose output already exists, and
writes the time each stage took to a JSON file beside the outputs. This is what makes
the run times in Section 3.3.2 measurements rather than estimates.

\subsection{Evaluation method}

Evaluation is described here once and is used throughout Chapter 5.

\textbf{Speech.} Word and character error rate are computed for each test clip
separately, and the median across clips is reported, because a single runaway output
can push a mean above 100 per cent. The base model and the adapted model are run on
exactly the same clips, so the comparison is paired and is tested with a Wilcoxon
signed-rank test and a sign test. The number of runaway clips is reported for both
models. Two spelling-tolerant variants are also computed. The first maps both sides
to the transcription guide's canonical spellings and collapses repeated vowels. The
second additionally counts a word as correct when it is at least 80 per cent similar
to the reference word, so that words of three letters or fewer still have to match
exactly. The rules for both were committed before any rescoring.

\textbf{Board reconstruction.} Two different things are counted and they must not be
confused. The fraction of tiles that found a completely clear view is computed by the
reconstruction itself and measures coverage, not whether writing survived. Whether
writing survived was answered by a person, who looked at each reconstructed board
beside two real frames from the same time and said whether anything written was
missing and what it was.

\textbf{Board reading and note quality.} Both are scored by recall of items from the
hand-verified answer keys. An answer key is a list of the things written on one
board, each labelled as a number, a name, a piece of code, a term or a phrase. Recall
is the right measure rather than F1, because good notes legitimately contain a great
deal that was not on the board, so precision would punish the behaviour we want.
Items of three characters or fewer must match as whole words, so that ``CS'' is not
found inside ``economics''. The keys were drafted from the board images and then
checked by hand, one board at a time, without the checker seeing any model output.
The folder holding them is refused by the note generator in code.

\textbf{Reporting rules.} Medians rather than means. The direction of every
difference is stated beside its p-value. Counting by board is reported alongside
counting by item, because the item count is dominated by a few dense boards and the
two can disagree. Results that go against the design are reported in the same detail
as results that support it.
